\documentclass[level=3]{kidscode}
\begin{document}
\kctitle{Notes for Families}{A letter about Level 3, privacy, and how to help at home}

\subsection{Dear families,}

Welcome to \textbf{Level 3: Advanced}, the last level of our home coding course. This letter is
for every family in the group, including friends who join us from another home. Please keep it
handy: the last page has a question to ask your child each week.

Over 12 weekly sessions of 60 minutes (plus two spare sessions for holidays or catching up), the
students will learn to work like real software makers: saving their work with \textbf{git},
designing programs with \textbf{classes and objects}, using famous \textbf{algorithms} for
searching, sorting and finding paths through mazes, exploring \textbf{data}, and building two
small \textbf{AI} programs to see how machines learn from examples, and how they can get things
wrong. The level ends with a \textbf{final project} of their own choice --- a grid game, a
generative art app or a data/AI app --- and a showcase you are warmly invited to.

Each week there is a short homework set (15--20 minutes) that never needs the internet.

\subsection{Our privacy promises}
\begin{safety}[What we promise, and why]
\begin{itemize}[leftmargin=1.2em,itemsep=2pt]
  \item \textbf{No accounts, no internet, no uploads.} Everything runs on our own Macs with free
        tools (Thonny, Python and Apple's git). Nothing your child makes is put online.
  \item \textbf{Nicknames only.} Git writes an author name and an email into every save point. We
        use a \textbf{nickname} and a made-up placeholder email that ends in
        \texttt{@example.com} (a web address reserved for examples, which never receives mail).
        Your child's real name, email address or school never goes into their projects.
  \item \textbf{No personal data in programs.} Programs ask for a nickname or a favourite colour,
        never a real name, age, address, school, photo or password. Survey data we analyse in
        week 8 is anonymous and collected only with permission.
  \item \textbf{AI chatbots are not needed.} No part of this course uses an AI chatbot. In week 9 we
        talk about using them safely (see below).
  \item \textbf{Showcase photos only with permission.} At the week 12 showcase we will only take
        photos or videos if every family agrees, and we will not share them online.
  \item \textbf{Projects stay yours.} Your child's project folder stays on the Mac it was made on.
        If you want a copy, we can put it on a USB stick for you.
\end{itemize}
\end{safety}

\subsection{AI chatbots at home: a few simple rules}
Many families now use AI chatbots. If yours does, these rules help keep children safe:
\begin{itemize}[itemsep=2pt]
  \item Never type real names, addresses, school names, phone numbers, passwords or photos into a
        chatbot.
  \item Check what it says. Chatbots can sound very sure and still be wrong. In week 9 your child
        builds a tiny \enquote{text generator} and sees why.
  \item Say when you used it, just like you would credit a friend who helped.
  \item Follow the service's age rules (many are 13+), and use it together with your child.
\end{itemize}

\subsection{How you can help}
\begin{itemize}[itemsep=2pt]
  \item \textbf{Ask your child to explain.} You do not need to know any coding. \enquote{Show me
        what it does} and \enquote{How did you make it do that?} are the best questions there are.
  \item \textbf{Do not fix their code.} When they are stuck, ask: \enquote{What did you expect?
        What happened instead? What could you try?} Getting unstuck is the skill we are building.
  \item \textbf{Celebrate bugs.} Every programmer makes mistakes all day. Finding one is a win.
  \item \textbf{Protect 15--20 minutes} for homework once a week, ideally soon after the session.
  \item \textbf{Mind screen time.} A stretch break every 20 minutes and the 20-20-20 eye rule
        (every 20 minutes, look 20 feet away for 20 seconds) help a lot.
\end{itemize}

\clearpage
\subsection{Ask your child about\ldots (one question per week)}

\renewcommand{\arraystretch}{1.45}
\begin{longtable}{@{}>{\bfseries}p{0.8cm}>{\raggedright\arraybackslash}p{4.4cm}p{\dimexpr\linewidth-5.2cm-4\tabcolsep}@{}}
\toprule
\rowcolor{kcaccent!15} Wk & This week & Ask\ldots\\ \midrule\endhead
\bottomrule\endfoot
1 & Save points with git & \enquote{What is a commit? Can you show me your list of save points?}\\ \midrule
2 & Classes and objects & \enquote{What is the difference between a blueprint and a thing made from it?}\\ \midrule
3 & Families of objects & \enquote{How is a Knight like a Character, and how is it different?}\\ \midrule
4 & Searching & \enquote{How could you find a word in a dictionary in only a few steps?}\\ \midrule
5 & Sorting, stacks and queues & \enquote{Is an undo button more like a pile of plates or a lunch line? Why?}\\ \midrule
6 & Review and checkpoint & \enquote{What bug did you find this week, and how did you track it down?}\\ \midrule
7 & Recursion and fractals & \enquote{How can a function call itself without going on forever?}\\ \midrule
8 & Data detectives & \enquote{What story did your chart tell? Was anything in the data messy?}\\ \midrule
9 & How machines learn & \enquote{How does your AI guess? When would it guess wrong?}\\ \midrule
10 & Project design & \enquote{What is your project, and which classes does it have?}\\ \midrule
11 & Testing and polishing & \enquote{What does one of your tests check? Did a test catch a bug?}\\ \midrule
12 & Showcase and graduation & \enquote{What design decision are you proudest of? What will you build next?}\\
\end{longtable}

\subsection{Questions or worries?}
Please talk to me at any time --- before or after a session is easiest.

\medskip
Teacher (parent) contact: \rule{8cm}{0.4pt}

\medskip
With thanks and excitement for the weeks ahead,

\medskip
\rule{6cm}{0.4pt}
\end{document}
